\documentclass[aspectratio=169,11pt]{beamer}
\usepackage{../../../shared/theme/esmad_beamer_theme}
\usepackage{amsmath,amssymb}
\usepackage{booktabs}
\usepackage{tikz}

\title[Causal Discovery]{Causal Discovery}
\subtitle{Assumptions, Graphs, and Limits}
\date{\today}

\begin{document}

\begin{frame}
\titlepage
\end{frame}

\begin{frame}{Learning Goals}
\begin{itemize}
\item distinguish causal discovery from causal effect estimation;
\item explain DAGs, d-separation, and Markov equivalence;
\item understand Markov and faithfulness assumptions;
\item compare constraint-based and score-based discovery;
\item explain how functional assumptions such as LiNGAM add orientation information;
\item recognize the role of latent confounding, interventions, and domain knowledge.
\end{itemize}
\end{frame}

\begin{frame}{Outline}
\tableofcontents
\end{frame}

\section{What Is Causal Discovery?}

\begin{frame}{The Structure-Learning Problem}
Suppose we observe variables
\[
X_1,\ldots,X_p.
\]

Causal discovery asks which directed structures are compatible with the observed statistical relationships and the maintained assumptions.

\begin{alertblock}{Different task}
Estimating a causal effect assumes a causal structure. Causal discovery tries to learn parts of that structure.
\end{alertblock}
\end{frame}

\begin{frame}{Graphs as Causal Models}
A directed acyclic graph (DAG) contains:
\begin{itemize}
\item nodes for variables;
\item directed edges for direct causal relations;
\item no directed cycles.
\end{itemize}

A graph encodes qualitative assumptions about the data-generating process.
\end{frame}

\section{Markov and Faithfulness}

\begin{frame}{Causal Markov Condition}
A DAG implies conditional independences through its graph structure.

Informally, each variable is independent of its non-descendants conditional on its parents.
\end{frame}

\begin{frame}{Factorization}
If a distribution is Markov with respect to DAG \(G\),
\[
p(x_1,\ldots,x_p)
=
\prod_{j=1}^{p}
p(x_j \mid \mathrm{pa}_G(x_j)).
\]

This links graph structure to statistical factorization.
\end{frame}

\begin{frame}{Faithfulness}
Faithfulness assumes that observed conditional independences arise from the graph structure rather than exact cancellations of causal effects.

\begin{alertblock}{Why it matters}
Without faithfulness, conditional-independence tests may remove or orient edges incorrectly.
\end{alertblock}
\end{frame}

\section{d-Separation}

\begin{frame}{Chains, Forks, and Colliders}
Three local patterns matter:

\[
X \rightarrow Z \rightarrow Y
\]

\[
X \leftarrow Z \rightarrow Y
\]

\[
X \rightarrow Z \leftarrow Y
\]

Chains and forks are blocked by conditioning on \(Z\).

Colliders behave differently: conditioning on \(Z\) can create dependence.
\end{frame}

\begin{frame}{Collider Bias}
For
\[
X \rightarrow Z \leftarrow Y,
\]
\(X\) and \(Y\) may be marginally independent.

Conditioning on \(Z\) can induce association.

\begin{block}{Connection}
This is why causal graphs help distinguish confounders from colliders.
\end{block}
\end{frame}

\section{Markov Equivalence}

\begin{frame}{Observational Data May Not Identify Direction}
Consider:
\[
X \rightarrow Y
\qquad\text{and}\qquad
X \leftarrow Y.
\]

With only two variables, purely observational data generally cannot distinguish these directions from conditional-independence structure alone.
\end{frame}

\begin{frame}{Equivalence Classes}
Two DAGs are Markov equivalent when they imply the same conditional independences.

They share:
\begin{itemize}
\item the same skeleton;
\item the same unshielded colliders.
\end{itemize}

Discovery often returns a partially directed graph rather than one unique DAG.
\end{frame}

\section{Constraint-Based Discovery}

\begin{frame}{PC Algorithm: Core Idea}
The PC algorithm:
\begin{enumerate}
\item starts from a dense undirected graph;
\item removes edges using conditional-independence tests;
\item identifies collider structures;
\item propagates edge orientations using logical rules.
\end{enumerate}
\end{frame}

\begin{frame}{Conditional-Independence Testing}
For Gaussian data, one may test partial correlations.

In other settings:
\begin{itemize}
\item kernel-based tests;
\item mutual-information tests;
\item nonparametric conditional-independence tests.
\end{itemize}

\begin{alertblock}{Finite-sample problem}
Discovery quality depends strongly on test power and calibration.
\end{alertblock}
\end{frame}

\begin{frame}{PC Assumptions}
Typical assumptions include:
\begin{itemize}
\item acyclicity;
\item causal Markov condition;
\item faithfulness;
\item no latent confounding in the basic version;
\item reliable conditional-independence tests.
\end{itemize}
\end{frame}

\section{Latent Confounding}

\begin{frame}{Hidden Common Causes}
Suppose an unobserved variable \(U\) causes both \(X\) and \(Y\):
\[
X \leftarrow U \rightarrow Y.
\]

A method that assumes causal sufficiency may infer a misleading direct edge between \(X\) and \(Y\).
\end{frame}

\begin{frame}{FCI-Type Methods}
FCI-style methods extend constraint-based discovery to allow latent confounding.

Outputs may contain:
\begin{itemize}
\item directed edges;
\item bidirected possibilities;
\item partially oriented endpoints.
\end{itemize}

The price is weaker orientation and more ambiguity.
\end{frame}

\section{Score-Based Discovery}

\begin{frame}{Search Over Graphs}
Score-based methods assign a score to candidate DAGs, such as:
\[
\mathrm{Score}(G)
=
\mathrm{fit}(G)
-
\mathrm{complexity\ penalty}(G).
\]

Then they search for a high-scoring graph.
\end{frame}

\begin{frame}{GES}
Greedy Equivalence Search operates over Markov equivalence classes.

It performs:
\begin{itemize}
\item a forward phase adding edges;
\item a backward phase removing edges.
\end{itemize}

Under strong assumptions and suitable scores, it can be consistent.
\end{frame}

\section{Functional Causal Models}

\begin{frame}{Beyond Conditional Independence}
Sometimes direction can be identified by assumptions on functional form and noise.

For example:
\[
Y=f(X)+\varepsilon,
\qquad
\varepsilon \perp X.
\]

The reverse model may not satisfy the same independence structure.
\end{frame}

\begin{frame}{LiNGAM}
LiNGAM assumes:
\begin{itemize}
\item linear relations;
\item acyclicity;
\item non-Gaussian independent disturbances.
\end{itemize}

These non-Gaussianity assumptions can identify causal ordering beyond ordinary covariance structure.
\end{frame}

\begin{frame}{Additive Noise Models}
In nonlinear settings, additive-noise asymmetry can sometimes distinguish
\[
X \rightarrow Y
\]
from
\[
X \leftarrow Y.
\]

Again, orientation comes from additional functional assumptions.
\end{frame}

\section{Interventions and Invariance}

\begin{frame}{Interventions Add Information}
Observational equivalence can be broken by interventions.

If we manipulate \(X\) and \(Y\) changes, that provides directional evidence that passive observation alone may not contain.
\end{frame}

\begin{frame}{Invariant Prediction}
Causal mechanisms are often assumed stable across environments.

If a conditional relation remains invariant while distributions shift, that stability can help identify plausible causal parents.
\end{frame}

\section{Time and Causal Direction}

\begin{frame}{Temporal Order Is Useful but Not Sufficient}
If \(X_t\) occurs before \(Y_{t+1}\), then reverse instantaneous causation is less plausible.

But time order alone does not remove:
\begin{itemize}
\item latent confounding;
\item common trends;
\item feedback through omitted variables;
\item measurement timing errors.
\end{itemize}
\end{frame}

\begin{frame}{Granger Causality Is Not Structural Causality}
If past \(X\) improves prediction of future \(Y\), then \(X\) Granger-causes \(Y\).

This is a predictive temporal concept, not automatically a structural causal claim.
\end{frame}

\section{Domain Knowledge}

\begin{frame}{Background Knowledge Can Constrain Search}
Useful constraints include:
\begin{itemize}
\item impossible directions;
\item known temporal ordering;
\item known exogenous variables;
\item forbidden edges;
\item required edges.
\end{itemize}

\begin{block}{Benefit}
Domain knowledge can shrink the equivalence class and improve stability.
\end{block}
\end{frame}

\section{Stability and Validation}

\begin{frame}{Bootstrap Graph Stability}
Repeat discovery across bootstrap samples.

Inspect:
\begin{itemize}
\item edge-selection frequency;
\item orientation frequency;
\item structural variability.
\end{itemize}

Unstable edges should not be presented as settled causal facts.
\end{frame}

\begin{frame}{Sensitivity to Tuning}
Discovery can change with:
\begin{itemize}
\item significance thresholds;
\item score penalties;
\item variable transformations;
\item included variables;
\item sample size.
\end{itemize}

Sensitivity is part of the result.
\end{frame}

\section{Applied Reasoning}

\begin{frame}{Example: Regional Pricing System}
Suppose we observe:
\begin{itemize}
\item regional price;
\item demand;
\item income;
\item competitor activity;
\item promotional intensity;
\item inventory constraints.
\end{itemize}

A discovery algorithm may suggest candidate structures, but economic knowledge is needed to rule out implausible directions and hidden-common-cause explanations.
\end{frame}

\begin{frame}{A Credible Causal-Discovery Workflow}
\begin{enumerate}
\item Define measured variables and temporal meaning.
\item State assumptions about latent confounding and cycles.
\item Encode domain constraints.
\item Choose a discovery family consistent with assumptions.
\item Run stability analysis.
\item Compare alternative algorithms.
\item Inspect the equivalence class, not only one graph.
\item Validate orientations using interventions or external evidence where possible.
\item Treat discovered structure as a hypothesis set, not automatic truth.
\end{enumerate}
\end{frame}

\begin{frame}{Common Failure Modes}
\begin{itemize}
\item claiming one observational DAG is uniquely identified when it is not;
\item ignoring latent confounding;
\item treating conditional independence tests as infallible;
\item using temporal prediction as proof of causality;
\item hiding unstable orientations;
\item forgetting that algorithm assumptions generate the conclusions.
\end{itemize}
\end{frame}

\begin{frame}{Synthesis: Discovery Is Assumption-Driven}
\begin{itemize}
\item DAGs encode qualitative causal structure.
\item Conditional independences can identify skeletons and some orientations.
\item Markov equivalence limits what passive observational data can reveal.
\item Faithfulness and causal sufficiency are strong assumptions.
\item Functional restrictions, non-Gaussianity, interventions, and invariance add identifying information.
\item Latent confounding and finite-sample error can radically alter conclusions.
\item Causal discovery should generate constrained hypotheses for further validation.
\end{itemize}
\end{frame}

\begin{frame}{Selected References}
\small
\begin{itemize}
\item Spirtes, P., Glymour, C., and Scheines, R. (2000). \textit{Causation, Prediction, and Search}, 2nd ed. MIT Press.
\item Pearl, J. (2009). \textit{Causality}, 2nd ed. Cambridge University Press.
\item Chickering, D. M. (2002). Optimal Structure Identification with Greedy Search. \textit{JMLR}, 3, 507--554.
\item Shimizu, S. et al. (2006). A Linear Non-Gaussian Acyclic Model for Causal Discovery. \textit{JMLR}, 7, 2003--2030.
\item Peters, J., Janzing, D., and Schölkopf, B. (2017). \textit{Elements of Causal Inference}. MIT Press.
\end{itemize}
\end{frame}

\contactslide

\end{document}
